%% notes-tex/wet-lab/hydrolysate-figure-gate/sections/3-remaining.tex
%% SOURCE: hand-authored status table. Each status repeats the word the panel script
%% draws; each "what remains" entry is a gap in the result files named under the figures
%% of sections/2-mockups.tex, or a round recorded as running or planned in
%% notes/experiments.040-inhibitor-synergy-wetlab.md (sections of 2026.10.10). Arm names
%% are the configuration groups of experiments/040-inhibitor-synergy-wetlab/conf/mixture/.

\clearpage
\section{What each panel still needs\secstatus{todo}}
\label{sec:remaining}

\begin{table}[htbp]
  \centering
  \footnotesize
  \caption[]{What each panel still needs before it is real. The status column repeats the
  panel's mark in words: on file (filled circle), every drawn value is in a committed
  result file; partial (half circle), a run is unfinished or a planned arm has no file;
  idea only (open circle), a layout or nothing run. No panel is approved; approval is the
  section status and is the author's.}
  \label{tab:remaining}
  \begin{tabular}{@{}p{22mm}p{7mm}p{141mm}@{}}
    \toprule
    panel & & what remains \\
    \midrule
    concept; main, a & \fstat{idea} &
      Redraw as one draw.io composition at print size. The counts in it are on file; the
      boxes and arrows are a layout, and the mixture path they describe (a set of
      compound tokens, each modulated by its own log-molar dose) is the proposed model
      and has not run. \\
    main, c & \fstat{idea} &
      Same: a layout whose numbers are on file. It merges into the concept sheet if the
      figure is cut from four rows to three. \\
    main, e; SI-1, a & \fstat{partial} &
      The second factorized round is producing. Two folds of the bilinear control and two
      of the environment encoder were on disk when the sheet was drawn; the round's
      remaining folds, the nine-seed ensembles and the stack with nested ridge are not
      written yet. Round eleven on Delta adds the width and depth axis. \\
    main, f & \fstat{partial} &
      The encoder arm of the learning curve has to be rerun on the corrected store; only
      the ridge arm has a file. The planned round also adds the other pooled sources to
      the gene-level task, which is the second line this panel would carry. \\
    main, i & \fstat{partial} &
      The transformer's compound-cold profiles for furfural and 5-HMF come out of the
      same running round. Until then the panel shows ridge against the measured profile
      and the screen's own repeat, and the transformer row is a label with no bar. \\
    main, j & \fstat{partial} &
      The model's predicted single-agent curves do not exist. They need the host head
      fitted on the titration run, in both variants the plan names: trained on the public
      dose anchors alone, and calibrated on the titration run. \\
    main, k, l & \fstat{partial} &
      The model's predictions over the combinations and over the isobole grids do not
      exist. The rules are on file and are the bars. If the model does not beat them, the
      panels ship with the rules and the model side by side and the text says so. \\
    main, b, d, g, h & \fstat{file} &
      Nothing. Every drawn value is in a committed result file. These four are what the
      figure can claim today. \\
    SI-1, b, c, d & \fstat{file} & Nothing. \\
    SI-2, all & \fstat{file} & Nothing. \\
    SI-3, all & \fstat{file} &
      Nothing, with one caveat the caption carries: the similarity matrices mix a
      measured profile for two inhibitors with predicted profiles for four, and the
      all-predicted version exists so that mixture is visible rather than hidden. \\
    \bottomrule
  \end{tabular}
\end{table}

\pfstep{The running and planned rounds, and what each one fills} The second factorized
round of the corrected-store campaign is on the two GilaHyper cards named for it, one job
for the bilinear control and one for the environment encoder; it fills main \textbf{e},
main \textbf{i} and the first column of SI-1 \textbf{a}. Round eleven of the same
campaign goes to Delta under the account the plan names, never a preemptible partition,
and widens the ladder rather than deepening it. The first mixture round of this
experiment is the one that fills the bottom row: about ten arms by three seeds, crossing
the pooled sources against the host-training recipe (the public dose anchors alone, the
anchors plus the titration run, and a fine-tune on the titration run) against the dose
treatment. It starts once the mixture trainer passes its processor smoke test, and its
launchers are the pair of scripts the experiment folder already carries.

\pfstep{The cut from four rows to three} The experiment note proposes merging rows two
and three. The merge that costs least is to move main \textbf{c} into the concept sheet,
where the same boxes already appear, and to put main \textbf{f} and main \textbf{g} side
by side as one panel on two axes, since both are about what more data buys: more
compounds on one axis, a second medium on the other. That leaves nine panels: the task,
the data, the ladder, the transformer against ridge, the merged data-and-dose panel, the
composition rule, the predicted profile, the single agents, and the combinations with the
isoboles as its inset. The alternative, dropping the single-agent panel, is worse: it is
what every rule in the bottom row is computed from, and a reader who cannot see the Hill
slopes cannot judge why Bliss and Loewe disagree as far as they do.
